
\documentclass[11pt]{article}
\usepackage[margin=1in]{geometry}
\usepackage{amsmath,amssymb,booktabs,hyperref,graphicx,parskip,enumitem}
\usepackage{xcolor}
\hypersetup{colorlinks=true,linkcolor=blue!50!black,urlcolor=blue!50!black}
\setlist{nosep}

\title{Replication Report: (title not recorded)\\[2pt]\large OSTI 1559043 \textbar{} Coverage 7/10, Agreement 7/10}
\author{Rick Stevens \and Ollie (AI Assistant)\\\small Argonne National Laboratory}
\date{2026-04-27}

\begin{document}\maketitle

\section{Source paper}
\begin{itemize}
\item \textbf{Title:} (title not recorded)
\item \textbf{Authors:} Not recorded
\item \textbf{Year:} n/a
\item \textbf{Domain:} Not recorded
\item \textbf{Identifier:} OSTI 1559043
\end{itemize}


\section{Replication objective}
The central claim of the paper concerns not recorded. Our objective was to reproduce the headline numerical/qualitative results within the constraints of the AI-driven Replication Project (24--72\,h, commodity GPU/CPU compute, no author contact). 

\section{Methodology}
We followed the standard Replication Project workflow: ingest the source PDF, extract method and data pointers, draft a replication plan, attempt the smallest end-to-end pipeline that produces a comparable number, then scale up where compute and data permit. Tools used in this replication are catalogued in the meta-paper's computational tools inventory (see \texttt{COMPUTATIONAL\_TOOLS\_INVENTORY.md}).

This paper went through the Tier-Lift pipeline; supplementary notes at \texttt{.openclaw/workspace/24h-progress/tier-lift-v2/1559043.md}.

\subsection*{Paper-directory README excerpt}
\small
\par\medskip\textbf{OSTI 1559043 — Jaravel et al.\textbackslash{} 2019 DNS of ignition-kernel stabilization in cross-flow}\par

**Replication status:** v4 (20-run Polaris PeleC ensemble) — **7/10**
(coverage 7/10, agreement 7/10). Promoted from v3 6/10.

\par\medskip\textbf{\# Latest report}\par
\par$\bullet$\ `replication-pelec/report\_v4/report.pdf`

\par\medskip\textbf{\# Ensemble summary (20-run Polaris, 2026-04-24)}\par

| φ    | N runs | N ignited | IP   | paper IP |
|------|--------|-----------|------|----------|
| 0.6  | 5      | 0         | 0.00 | 0.00     |
| 0.8  | 5      | 0         | 0.00 | 0.20     |
| 1.0  | 5      | 5         | 1.00 | 0.65     |
| 1.2  | 5      | 5         | 1.00 | 0.90     |

Ignition criterion: `T\_max(0.9 ms) \textgreater{} 2550 K` for complete runs **or**
`max T\_max(t\textgreater{}0.5 ms) \textgreater{} 2700 K` for CFL-truncated runs. See
`replication-pelec/report\_v4/report.tex` §3.

Per-run details: `replication-pelec/polaris\_ensemble/summary.json`.

\par\medskip\textbf{\# Key corrections vs prior notes}\par
\par$\bullet$\ Aurora job **8449654 never ran** (reservation-blocked). Only 8449560
  ran — single φ=0.6 realisation, 9 checkpoints. Aurora is **not** the
  ensemble source.
\par$\bullet$\ Polaris job set 7099651–7099670 (20 jobs) is the real ensemble.

\par\medskip\textbf{\# Directory map}\par

[code block omitted]

\par\medskip\textbf{\# Known limitations (see report\_v4 §5)}\par

1. 128×64×64 grid, no AMR level 1 — turbulence under-resolved
2. 1-ms simulation window (paper goes to \textasciitilde{}10 ms)
3. Synthetic (not paper's tabulated) turbulent inflow
4. 3 of 5 φ=1.2 runs CFL-truncated at 0.25–0.63 ms; classified ignited
   from elevated late-window T\_max
5. Per-plotfile species-max extraction on Polaris debug queue; at report
   time still queued. Hook in `polaris\_ensemble/analyze.py` consumes the
   per-run `*\_timeseries.csv` when available to add species figures.
\normalsize

The replication was driven by an OpenClaw agent loop (Claude Opus / GPT-5 class models) issuing shell, Python, and (where applicable) GPU jobs. Compute usage and wall-time for individual papers are summarised in the meta-paper's resource table; not separately recorded for this paper unless noted in the Tier-Lift artefact. Sub-agents were spawned for replication-plan generation and (for papers in batches 1--4) for tier-lift extension runs.

\section{What was reproduced}
\begin{itemize}
\item Not recorded.
\end{itemize}

\section{What was missing or incomplete}
\begin{itemize}
\item Not recorded.
\end{itemize}
The dominant blocker categories for this paper, mapped onto the meta-paper's 9-category friction taxonomy (\S4), are listed in \S\ref{sec:friction} below.

\section{Quantitative agreement}
Per-quantity numerical comparisons are not separately tabulated in the structured evaluation record for this paper beyond the agreement rationale below; where the rationale references specific numbers, they are reproduced verbatim. A full numerical comparison table is recorded in the per-replication artefacts (see \S\ref{sec:repro}). For papers below 8/8, the agreement rationale identifies the specific quantities that diverge.

\paragraph{Agreement rationale.} Not recorded.

\section{Score and friction-point tagging}\label{sec:friction}
\begin{itemize}
\item \textbf{Coverage:} 7/10 --- Not recorded.
\item \textbf{Agreement:} 7/10 --- Not recorded.
\end{itemize}

\noindent\textbf{Friction points encountered (taxonomy tags):}
\begin{itemize}
\item \textbf{F5}: Force-field / hyperparameter drift (unspecified configuration)
\end{itemize}

\section{Follow-on questions}
\begin{itemize}
\item Not recorded.
\end{itemize}

\section{Reproducibility pointers}\label{sec:repro}
\begin{itemize}
\item \textbf{Paper directory:} \texttt{1559043-ignition-kernel-turbulent}
\item \textbf{Replication artefacts:}
\begin{itemize}
\item \texttt{/Users/stevens/Dropbox/REPLICATE-PROJECT/1559043-ignition-kernel-turbulent/replication}
\end{itemize}
\item \textbf{Replication-dir hint from JSONL:} \texttt{/Users/stevens/Dropbox/REPLICATE-PROJECT/1559043-ignition-kernel-turbulent}
\item \textbf{Status:} tier\_lift\_v2\_2026\_04\_27
\end{itemize}

To re-run, change directory into the paper directory above and consult its \texttt{README.md} for the entry-point script. Most replications follow the convention \texttt{replication/run\_*.sh} or \texttt{replication/<subdir>/run.py}.

\end{document}
